\subsection{The $2\times2$ Case}

Let
\[
A=\begin{pmatrix}a&b\\c&d\end{pmatrix}.
\]
Then
\[
\det A=0\iff ad=bc.
\]

\begin{proofbox}[Zero determinant and proportional rows]
Assume first that the first row is nonzero. If $a\neq0$, then
\[
d=\frac{bc}{a},
\qquad
(c,d)=\frac ca(a,b).
\]
If $a=0$, then $b\neq0$ and $ad=bc$ gives $c=0$, so
\[
(c,d)=\frac db(0,b).
\]
If the first row is zero, the relation is immediate.
\end{proofbox}

\begin{theorembox}[Dependence criterion in dimension two]
For a $2\times2$ matrix,
\[
\det A=0
\]
if and only if its rows are proportional. The same statement holds for its columns.
\end{theorembox}
